% !TeX program = lualatex
% !TeX encoding = utf8
% !TeX spellcheck = uk_UA
% !TeX root = ../EMConspect.tex

%=========================================================
\chapter{Магнітне поле в речовині}\label{\currfilebase}
\graphicspath{{\currfilebase/Pictures/}}
%=========================================================
\epigraph{Речовина не створює нового типу поля --- вона лише ховає у собі мільярди мільярдів мікроскопічних струмів, і все мистецтво теорії
	магнетизму в речовині полягає в тому, щоб навчитися не рахувати їх поштучно.
}{Переосмислення ідей сучасної фізики}

\section*{Гіпотеза Ампера}

Ми вже знаємо, що магнітне поле діє на рухомі заряджені частинки та на рамки зі струмом. Але \emphx{чому воно діє також і на шматок
	звичайного заліза} --- скажімо, на цвях, який сам по собі не є магнітом? Відповідь на це запитання дав ще А.~Ампер: він припустив,
	що \emphx{всередині магніту теж течуть струми}.

Якщо взяти магнітну стрілку чи магніт у руки, жодних струмів ми, звісно, не відчуваємо: отже, ці струми циркулюють \emphx{усередині} речовини і
ніколи не виходять назовні. Що це за струми, Ампер, звісно, не знав --- адже про будову атома тоді не було відомо нічого. Сучасній науці вже
відомо, що звичайна речовина складається з атомів: усередині атомів є позитивно заряджені ядра, навколо яких рухаються негативно заряджені
електрони, причому самі електрони, до того ж, --- маленькі магнітні \enquote{стрілочки} (мають власний магнітний момент, спін). Отже, рух
електронів усередині атомів --- це і є ті самі \emphx{струми}, про існування яких здогадався Ампер; магнітне поле діє на ці струми, а отже --- і
на речовину загалом.

\begin{figure}[h!]\centering
	\includegraphics[width=0.6\linewidth]{AmpereHypotesis}
	\caption{Гіпотеза Ампера: магнетизм речовини зумовлений мікроскопічними
		(молекулярними) струмами всередині неї\label{pic:AmpereHypotesis}}
\end{figure}

%% --------------------------------------------------------
\section{Основні поняття}
%% --------------------------------------------------------

\subsection*{Мікрополе та середнє поле}

У речовині магнітне поле формується як зовнішнім полем, так і струмами, що циркулюють у самій речовині. На мікрорівні (тобто на відстанях порядку
розміру атома й менше) це поле різко змінюється в часі й просторі --- таке поле називають \emphx{мікрополем} $\Bfield_\text{micro}$. Однак якщо
провести усереднення за \emphx{фізично нескінченно малим} об'ємом $\Delta V$ (малим порівняно з масштабом, на якому поле помітно змінюється
макроскопічно, але великим порівняно з розмірами й відстанями між окремими атомами --- тобто таким, що вміщує величезну кількість атомів), отримаємо
\emphx{середнє поле}:
\begin{equation}\label{eq:AvgField}
	\tcbhighmath{
		\langle\Bfield\rangle = \frac1{\Delta V}\iiint\limits_{\Delta V} \Bfield_\text{micro}\, dV.
	}
\end{equation}
Середнє поле змінюється значно повільніше за мікрополе завдяки статистичному усередненню за хаотичного теплового руху частинок. Надалі, як і
раніше з полем $\Efield$ в діелектрику, під $\Bfield$ ми завжди розумітимемо саме це \emphx{середнє} поле.

\begin{Attention}
	Ця конструкція --- точна аналогія переходу від мікрополя до макроскопічного поля $\Efield$, яку ми обговорювали в розділі про поле в
	діелектриках (гл.~\ref{FileldInDielectrics}): у макроскопічній електродинаміці ми принципово ніколи не працюємо з полем окремого атома, а лише
	з полем, усередненим за фізично нескінченно малим об'ємом. Цей самий прийом --- спільна методологічна основа всієї макроскопічної теорії поля
	в речовині, як електричній, так і магнітній.
\end{Attention}

\subsection*{Струми провідності та молекулярні струми}

Усі струми, створювані рухомими зарядами в речовині, поділяють на дві принципово різні групи:
\begin{enumerate}
	\item \emphx{Струми провідності} --- пов'язані з переміщенням \emphx{вільних} носіїв заряду (тих самих, що ми вивчали в
	      гл.~\ref{SteadyCurrent}); це \emphx{макроскопічні} струми, зовнішні щодо самої речовини й підтримувані стороннім джерелом (ЕРС).
	\item \emphx{Молекулярні струми} --- зумовлені орбітальним рухом електронів в атомах (молекулах) та їхнім спіном (власним моментом імпульсу), а
	      також, значно меншою мірою, магнітними моментами ядер. Це \emphx{внутрішні}, властиві самій речовині струми.
\end{enumerate}

%% --------------------------------------------------------
\section{Намагнічування}
%% --------------------------------------------------------

\subsection*{Вектор намагніченості}

\emphx{Вектором намагніченості} (\emphx{намагніченістю}) називають величину, що характеризує сумарний магнітний момент одиниці об'єму речовини:
\begin{equation}\label{eq:MagnetizationDef}
	\tcbhighmath{
		\vect J = \frac1{\Delta V}\sum\limits_{i\,\in\,\Delta V} \vect p_i,
	}
	\qquad \vect p_i \text{ --- магнітний момент $i$-ї частинки,}
\end{equation}
де $\Delta V$ --- той самий \emphx{фізично нескінченно малий} об'єм, що й у означенні~\eqref{eq:AvgField}.
Намагніченість називають \emphx{однорідною}, якщо вектор $\vect J$ не залежить від точки в речовині; якщо $\vect J \neq \const$ --- \emphx{неоднорідною}.

\begin{Attention}
	Вектор намагніченості $\vect J$ --- точний магнітний аналог вектора поляризації $\vect P$ в діелектриках (гл.~\ref{FileldInDielectrics}): там
	$\vect P$ --- густина електричних дипольних моментів, тут $\vect J$ --- густина магнітних дипольних моментів. Уся подальша логіка цього
	розділу (виведення зв'язаних струмів, введення допоміжного вектора $\Hfield$, граничні умови) дослівно повторює логіку розділу про
	діелектрики --- з тією різницею, що диполі там \emphx{електричні}, а тут --- \emphx{магнітні}.
\end{Attention}

\subsection*{Основна задача теорії магнітостатики в речовині}

Пряме обчислення поля в намагніченій речовині вимагало б підсумовування магнітних полів від колосальної кількості (порядку числа Авогадро)
молекулярних струмів --- завдання абсолютно нездійсненне практично. Основна ідея теорії магнітостатики в речовині --- \emphx{вилучити} молекулярні
струми з рівнянь, замінивши їх іншою, макроскопічно \emphx{вимірюваною} величиною --- вектором намагніченості $\vect J$. Це, за аналогією з тим, що
ми вже робили для зв'язаних зарядів у діелектриках, істотно спрощує розрахунок полів у магнетиках.

\subsection*{Зв'язок намагніченості з молекулярними струмами}

Виділимо в намагніченій речовині малий косокутний циліндр, орієнтований так, щоб поле в ньому можна було вважати практично однорідним
(рис.~\ref{pic:Cylinder}). У товщі циліндра кільцеві молекулярні струми сусідніх атомів течуть у протилежних напрямках і взаємно компенсуються --- в
усіх точках, \emphx{окрім} точок бічної поверхні (рис.~\ref{pic:CylinderTop}, вигляд торця): там компенсація неповна, і залишається нескомпенсований
поверхневий струм, що тече бічною поверхнею циліндра.

\begin{figure}[h!]
	\centering
	\begin{minipage}[t]{0.48\linewidth}\centering
		\localinput{cylinder.tikz}
		\caption{Косокутний циліндр у намагніченій речовині}\label{pic:Cylinder}
	\end{minipage}\hfill
	\begin{minipage}[t]{0.48\linewidth}\centering
		\localinput{cylindertop.tikz}
		\caption{Вигляд торця циліндра: кільцеві струми компенсуються всюди, крім бічної поверхні}\label{pic:CylinderTop}
	\end{minipage}
\end{figure}

Магнітний момент такого циліндрика, з одного боку, дорівнює $\vect p_m = \vect J V$ (за означенням намагніченості), а з іншого --- визначається
нескомпенсованим поверхневим струмом $I_m$, що тече по бічній поверхні, за формулою для магнітного моменту витка~\eqref{eq:MagMomentLoop}
(гл.~\ref{MagFieldVacuum}), $\vect p_m = \frac1c I_m S\vect n$, де $S$ --- площа торця циліндра, $\vect n$ --- нормаль до торця, а $\theta$ ---
кут між твірною $\vect\tau$ циліндра і цією нормаллю. Об'єм косокутного циліндра дорівнює $V = S\,\ell\cos\theta$, де $\ell$ --- довжина
твірної. Прирівнюючи обидва вирази й проєктуючи рівність на напрям твірної:
\begin{equation*}
	\vect J\, V = \frac1c I_m S\vect n \quad\Rightarrow\quad (\vect J\cdot\vect\tau)\, S\,\ell\cos\theta = \frac1c I_m\, S\,(\vect n\cdot\vect\tau),
	\qquad \vect n\cdot\vect\tau = \cos\theta.
\end{equation*}
Скорочуючи на $S\cos\theta$ і ділячи на $\ell$, знаходимо \emphx{поверхневу густину молекулярного струму} $i_m = I_m/\ell$:
\begin{equation}\label{eq:MolecularSurfaceCurrent}
	\tcbhighmath{
		i_m = c\, \vect J\cdot\vect\tau.
	}
\end{equation}

\begin{Attention}
	Формула~\eqref{eq:MolecularSurfaceCurrent} --- скалярний запис векторного співвідношення $\vect i_m = c\,\vect J\times\vect n$, з якого
	видно, що \emphx{поверхневий} молекулярний струм перпендикулярний до намагніченості: $\vect i_m \perp \vect J$. Це разюче
	відрізняється від електричного аналога: там зв'язаний поверхневий заряд $\sigma' = P_n$ визначався \emphx{нормальною} складовою $\vect P$ (гл.~\ref{FileldInDielectrics}), тут же зв'язаний струм визначається \emphx{тангенціальною} складовою $\vect J$. Причина
	--- принципова: $\vect P$ і заряд пов'язані через дивергенцію (скалярне джерело), тоді як $\vect J$ і струм --- через ротор (векторний,
	вихровий об'єкт); цю відмінність ми зараз узагальнимо на довільний контур.
\end{Attention}

%% --------------------------------------------------------
\section{Теорема про циркуляцію в речовині}
%% --------------------------------------------------------

\subsection*{Циркуляція вектора намагніченості}

Виберемо тепер у речовині \emphx{довільний} замкнений контур $L$ (рис.~\ref{pic:Jloop}). За формулою~\eqref{eq:MolecularSurfaceCurrent} на елемент довжини
контуру $d\vect\ell$ припадає молекулярний струм $dI_m = i_m\, d\ell = c\,\vect J\cdot d\vect\ell$, тому повний молекулярний струм, що
перетинає (пронизує) контур,
\begin{equation}\label{eq:CirculationJ}
	I_m = c\oint\limits_L \vect J\cdot d\vect\ell.
\end{equation}

\begin{SCfigure}[2][h!]\centering
	\localinput{Jloop.tikz}
	\caption{Довільний замкнений контур $L$ у намагніченій речовині\label{pic:Jloop}}
\end{SCfigure}

За теоремою Стокса цей вираз перетворюється так:
\begin{equation*}
	I_m = c\oint\limits_L \vect J\cdot d\vect\ell = c\iint\limits_S \rot\vect J\cdot d\vect S,
\end{equation*}
де $S$ --- будь-яка поверхня, натягнута на контур $L$. З іншого боку, той самий струм $I_m$ можна виразити через його \emphx{об'ємну} густину
$\vect j_m$: $I_m = \iint_S \vect j_m\cdot d\vect S$. Порівнюючи два вирази, отримуємо \emphx{зв'язок молекулярного струму з вектором
	намагніченості в диференціальній формі}:
\begin{equation}\label{eq:MolecularCurrentDensity}
	\tcbhighmath{
		\vect j_m = c\rot\vect J.
	}
\end{equation}

\subsection*{Вектор напруженості магнітного поля}

Циркуляцію повного магнітного поля $\Bfield$ у речовині породжують \emphx{усі} струми --- і струми провідності $I$, і молекулярні струми $I_m$:
\begin{equation*}
	\oint\limits_L \Bfield\cdot d\vect r = \frac{4\pi}c\left(I + I_m\right) = \frac{4\pi}c\left(I + c\iint\limits_S \rot\vect J\cdot d\vect
	S\right).
\end{equation*}
Щоб \emphx{вилучити} невимірюваний молекулярний струм $I_m$ з цього рівняння, введемо допоміжну величину --- \emphx{напруженість магнітного поля}:
\begin{equation}\label{eq:HDef}
	\tcbhighmath{
		\Hfield = \Bfield - 4\pi\vect J.
	}
\end{equation}
Підставляючи це означення в попереднє рівняння, молекулярний струм скорочується, і \emphx{теорема про циркуляцію в речовині} набуває вигляду
\begin{equation}\label{eq:CirculationHInt}
	\tcbhighmath{
		\oint\limits_L \Hfield\cdot d\vect r = \frac{4\pi}c I.
	}
\end{equation}
Еквівалентна диференціальна форма цієї теореми:
\begin{equation}\label{eq:CirculationHDiff}
	\tcbhighmath{
		\rot\Hfield = \frac{4\pi}c\vect j.
	}
\end{equation}

\begin{Attention}
	Вектор $\Hfield$ --- \emphx{допоміжна} (не фундаментальна) величина, що слугує лише для спрощення вигляду рівнянь; принципово важливо, що
	\emphx{його циркуляція визначається тільки струмами провідності} --- тими, які ми задаємо чи вимірюємо ззовні, а не невідомими молекулярними
	струмами. Це прямий магнітний аналог того, як вектор $\Dfield = \Efield + 4\pi\vect P$ в електростатиці дав змогу виключити невідомі зв'язані
	заряди й записати теорему Гаусса лише через вільні заряди (гл.~\ref{FileldInDielectrics}). Аналогія майже дослівна --- лише $\Dfield$ виникає з
	теореми \emphx{Гаусса} (через дивергенцію $\vect P$), а $\Hfield$ --- із теореми про \emphx{циркуляцію} (через ротор $\vect J$), що знову
	відображає різну (потенціальну проти вихрової) природу цих двох полів.
\end{Attention}

%% --------------------------------------------------------
\section{Магнітна сприйнятливість речовини}
%% --------------------------------------------------------

%% --------------------------------------------------------
\subsection*{Лінійні ізотропні середовища}
%% --------------------------------------------------------

Розглянемо \emphx{лінійні ізотропні} середовища --- ті, в яких
намагніченість пропорційна полю, а магнітні властивості однакові в усіх
напрямках. Середовище (речовина), що реагує на зовнішнє магнітне поле, називатимемо
\emphx{магнетиком}. Якщо поле не надто сильне, а в середовищі немає
початкової (залишкової) намагніченості, намагніченість, як показує дослід,
прямо пропорційна полю:
\begin{equation}\label{eq:LinearMedium}
	\vect J = \chi_m\Hfield,
\end{equation}
де $\chi_m$ --- \emphx{магнітна сприйнятливість} середовища (безрозмірна
величина). Підставляючи це в означення~\eqref{eq:HDef},
$\Bfield = \Hfield + 4\pi\vect J = (1+4\pi\chi_m)\Hfield$, тобто
\begin{equation}\label{eq:BMuH}
	\tcbhighmath{
		\Bfield = \mu\Hfield, \qquad \mu = 1+4\pi\chi_m,
	}
\end{equation}
де величину $\mu$ називають \emphx{магнітною проникністю середовища}.

%% --------------------------------------------------------
\subsection*{Класифікація магнетиків}
%% --------------------------------------------------------

Магнетиком є будь-яка речовина: адже, на відміну від діелектриків,
молекулярні струми --- а отже, і магнітні моменти атомів --- є в кожному
без винятку атомі: рух електронів навколо ядра не припиняється ніколи.
Різні речовини відрізняються не наявністю чи відсутністю цих мікрострумів,
а тим, як вони поводяться в зовнішньому полі і чи компенсують вони одне
одного за його відсутності.

Залежно від знака й величини магнітної сприйнятливості виділяють такі
основні класи лінійних магнетиків:
\begin{enumerate}
	\item Якщо $\chi_m > 0$ (а отже, $\mu>1$), магнетик називають
	      \emphx{парамагнетиком}. Парамагнітні властивості мають, наприклад,
	      \ce{Al}, \ce{Pt}, \ce{FeCl2}, \ce{O2}, лужні й лужноземельні метали.
	\item Якщо $\chi_m < 0$ (а отже, $\mu<1$), магнетик називають
	      \emphx{діамагнетиком}. Діамагнетиками є \ce{Bi}, \ce{Sb}, \ce{Si},
	      \ce{H2O}, \ce{H2}, \ce{N2} тощо.
\end{enumerate}

І парамагнетики, і діамагнетики належать до класу лінійних ізотропних
середовищ, розглянутих вище.

Деякі речовини здатні зберігати намагніченість $\vect J$ навіть за
\emphx{відсутності} зовнішнього поля --- для них лінійне
співвідношення~\eqref{eq:LinearMedium} узагалі не виконується. Такі
речовини називають \emphx{феромагнетиками} (\ce{Fe}, \ce{Co}, \ce{Ni} та
багато їхніх сплавів). У тому діапазоні полів, де лінійне наближення все ж
формально можна застосувати, магнітна проникність феромагнетиків сягає
величезних значень, $\mu \gg 1$ (тисячі й навіть десятки тисяч) ---
детальніше про них ми поговоримо в останньому розділі цієї глави.

%\begin{Attention}
%	Фізична природа діа- і парамагнетизму --- принципово різна. \emphx{Діамагнетизм} --- явище \emphx{універсальне}: зовнішнє поле, за законом
%	електромагнітної індукції (детальніше --- у відповідному розділі курсу), незначно змінює швидкість обертання електронів на орбіталях
%	(\emphx{прецесія Лармора}), наводячи додатковий молекулярний струм, напрямлений \emphx{проти} поля (за правилом Ленца) --- тому діамагнітний
%	внесок завжди від'ємний і присутній \emphx{в усіх без винятку} речовинах. \emphx{Парамагнетизм} виникає лише тоді, коли атоми (молекули)
%	речовини мають власний (нескомпенсований) магнітний момент навіть без поля: тоді зовнішнє поле частково \emphx{орієнтує} ці вже наявні моменти
%	вздовж себе (борючись із тепловим хаотичним рухом, що цьому протидіє), даючи додатний внесок у $\chi_m$, який, коли він є, зазвичай на порядки
%	переважає слабший діамагнітний внесок. Тому речовина виявляється парамагнітною, якщо в її атомах є нескомпенсований момент, і діамагнітною ---
%	якщо всі моменти в основному стані скомпенсовані.
%\end{Attention}




%=========================================================
\begin{problem}%
Знайдемо індукцію магнітного поля всередині нескінченного соленоїда з густиною намотки $n$ витків на одиницю довжини, по витках якого тече струм
$I$, якщо об'єм усередині соленоїда заповнено парамагнетиком з проникністю $\mu$.
\end{problem}

%=========================================================
\begin{problem}%
Постійний струм $I$ тече вздовж довгого циліндричного провідника круглого перерізу радіуса $R$, виготовленого з парамагнетика зі сприйнятливістю
$\chi$ (рис.~\ref{pic:CurrentInMagnetic}). Знайдемо поверхневий і об'ємний молекулярні струми та з'ясуємо, як вони напрямлені один відносно
одного.
\end{problem}



%
%
%
%
%\subsection*{Приклад 2: молекулярні струми в парамагнітному провіднику}
%
%Постійний струм $I$ тече вздовж довгого циліндричного провідника круглого перерізу радіуса $R$, виготовленого з парамагнетика зі сприйнятливістю
%$\chi$ (рис.~\ref{pic:CurrentInMagnetic}). Знайдемо поверхневий і об'ємний молекулярні струми та з'ясуємо, як вони напрямлені один відносно
%одного.
%
%\begin{SCfigure}[3][h!]\centering
%	\localinput{current_in_magnetic.tikz}
%	\caption{Циліндричний провідник з парамагнетика зі струмом $I$: молекулярні
%		струми (умовно показані малими витками) течуть навколо кожного елемента об'єму\label{pic:CurrentInMagnetic}}
%\end{SCfigure}
%
%Поле $\Hfield$ усередині провідника (на відстані $r<R$ від осі), як і в звичайному провіднику без магнетика (гл.~\ref{MagFieldVacuum}, формула
%типу~\eqref{eq:CylinderInside}, тільки для $\Hfield$), знаходимо за теоремою про циркуляцію, охоплюючи струм $I(r) = I r^2/R^2$: $H\cdot 2\pi r =
%\frac{4\pi}c I \dfrac{r^2}{R^2}$, тобто $\Hfield$ напрямлене по дотичній до кіл, співвісних із провідником (азимутально), і
%\begin{equation*}
%	H(r) = \frac{2Ir}{cR^2}, \qquad J(r) = \chi H(r) = \frac{2\chi I r}{cR^2} \qquad (\text{також азимутально}).
%\end{equation*}
%
%\emphx{Об'ємний молекулярний струм.} Оскільки $\vect J$ має лише азимутальну складову $J_\varphi(r)$, за формулою для ротора в циліндричних
%координатах $(\rot\vect J)_z = \dfrac1r\dfrac{d(rJ_\varphi)}{dr}$, тому за формулою~\eqref{eq:MolecularCurrentDensity}
%\begin{equation*}
%	j_{m,z} = \frac{c}r\frac{d}{dr}\left(r\cdot\frac{2\chi Ir}{cR^2}\right) = \frac{c}r\cdot\frac{4\chi Ir}{cR^2} = \frac{4\chi I}{R^2}.
%\end{equation*}
%Об'ємний молекулярний струм однорідно розподілений по перерізу й тече \emphx{вздовж осі провідника, у тому самому напрямку}, що й струм
%провідності $I$; повний об'ємний молекулярний струм
%\begin{equation}\label{eq:WireVolCurrent}
%	\tcbhighmath{
%		I'_\text{об} = j_{m,z}\,\pi R^2 = 4\pi\chi I.
%	}
%\end{equation}
%
%\emphx{Поверхневий молекулярний струм.} На поверхні провідника ($r=R$) намагніченість $\vect J$ стрибком спадає до нуля (зовні провідника
%магнетика немає), що, за формулою~\eqref{eq:MolecularSurfaceCurrent} (з $\vect\tau = \vect\varphi\times\vect n$, де $\vect n$ --- зовнішня
%нормаль), дає поверхневий молекулярний струм, напрямлений \emphx{уздовж осі провідника, але у протилежному до $I$ напрямку}, з густиною $i_m =
%cJ(R) = 2\chi I/R$; повний поверхневий молекулярний струм
%\begin{equation}\label{eq:WireSurfCurrent}
%	\tcbhighmath{
%		I'_\text{пов} = i_m\cdot 2\pi R = 4\pi\chi I.
%	}
%\end{equation}
%
%\begin{Attention}
%	Обидва молекулярні струми мають \emphx{однакову за величиною} силу $4\pi\chi I$, але напрямлені \emphx{протилежно}: об'ємний --- уздовж
%	струму провідності, поверхневий --- проти нього. Це не випадковість, а наслідок закону збереження заряду: молекулярні струми --- це
%	\emphx{замкнені} струми (вони не переносять заряд крізь переріз провідника в цілому, на відміну від струму провідності), тому сумарний
%	молекулярний струм крізь будь-який переріз циліндра мусить дорівнювати нулю, $I'_\text{об} - I'_\text{пов} = 0$ --- рівність, яку ми якраз і
%	отримали прямим обчисленням.
%\end{Attention}

%% --------------------------------------------------------
\section{Граничні умови}
%% --------------------------------------------------------

\subsection*{Перша гранична умова (для нормальної складової $\Bfield$)}

Застосуємо теорему Гаусса для вектора $\Bfield$ (яка, нагадаємо, залишається справедливою в будь-якому середовищі, оскільки джерел магнітного поля,
на відміну від електричного, --- магнітних зарядів --- не існує) до нескінченно малого циліндра, що охоплює невелику ділянку межі розділу
двох магнетиків (рис.~\ref{pic:BC1Mag}). Вважаючи $d\vect S_2 = -d\vect S_1 = \vect n\, dS$ ($\vect n$ --- нормаль від середовища 1 до
середовища 2, як на рисунку):
\begin{equation*}
	\oiint\limits_S \Bfield\cdot d\vect S = 0 \quad\Rightarrow\quad \Bfield_1\cdot d\vect S_1 + \Bfield_2\cdot d\vect S_2 = 0.
\end{equation*}

\begin{figure}[h!]\centering
	\localinput{bc1.tikz}
	\caption{Циліндр Гаусса на межі розділу двох магнетиків\label{pic:BC1Mag}}
\end{figure}

Звідси перша гранична умова:
\begin{equation}\label{eq:BC1Magnetic}
	\tcbhighmath{
		B_{1n} = B_{2n}
	}
\end{equation}
--- \emphx{нормальна складова вектора $\Bfield$ не зазнає стрибка} при переході через межу розділу середовищ (дослівний магнітний аналог
неперервності $D_n$ за відсутності вільного заряду на межі, гл.~\ref{FileldInDielectrics}).

\subsection*{Друга гранична умова (для тангенціальної складової $\Hfield$)}

Застосуємо теорему про циркуляцію~\eqref{eq:CirculationHInt} до нескінченно малого прямокутного контуру $ABCD$, довгі сторони якого проходять
на нескінченно малій відстані по обидва боки від межі розділу (рис.~\ref{pic:BC2Mag}). Поверхневий струм провідності (якщо він там є)
характеризуємо одиничним вектором $\vect b$ уздовж напряму струму й лінійною густиною $i_b$; контур орієнтуємо \emphx{впоперек} струму, так щоб
струм пронизував його, тобто вздовж орта $\vect\tau = \vect b\times\vect n$. Обходячи контур так, що $d\vect r_2 = \vect\tau\, d\ell$ (сторона в
середовищі 2) і $d\vect r_1 = -d\vect r_2$:
\begin{equation*}
	\oint\limits_L \Hfield\cdot d\vect r = \frac{4\pi}c I \quad\Rightarrow\quad \Hfield_1\cdot d\vect r_1 + \Hfield_2\cdot d\vect r_2 =
	\frac{4\pi}c\, i_b\, d\ell.
\end{equation*}

\begin{figure}[h!]\centering
	\localinput{bc2.tikz}
	\caption{Контур циркуляції на межі розділу двох магнетиків\label{pic:BC2Mag}}
\end{figure}

Звідси друга гранична умова:
\begin{equation}\label{eq:BC2MagneticScalar}
	\tcbhighmath{
		H_{2\tau} - H_{1\tau} = \frac{4\pi}c\, i_b.
	}
\end{equation}
У векторному вигляді, циклічно переставивши співмножники в мішаному добутку, цю умову записують як
\begin{equation}\label{eq:BC2MagneticVector}
	\tcbhighmath{
		\vect n\times\left(\Hfield_2 - \Hfield_1\right) = \frac{4\pi}c\, \vect i_b.
	}
\end{equation}

\begin{Attention}
	Тангенціальна складова $\Hfield$ зазнає стрибка \emphx{тільки} там, де по поверхні розділу середовищ течуть \emphx{струми провідності} $i_b$
	(наприклад, обмотка, намотана безпосередньо на межі двох магнетиків); якщо на межі струмів провідності немає, $H_{1\tau}=H_{2\tau}$ ---
	тангенціальна складова $\Hfield$ неперервна, а от тангенціальна складова $\Bfield = \mu\Hfield$, узагалі кажучи, \emphx{стрибає} (оскільки
	$\mu_1\neq\mu_2$) --- цілком аналогічно до того, як на межі діелектриків без вільного заряду неперервна $D_n$, а не $E_n$
	(гл.~\ref{FileldInDielectrics}).
\end{Attention}

\subsection*{Заломлення силових ліній на межі магнетиків}

Розглянемо межу розділу двох магнетиків з проникностями $\mu_1$ і $\mu_2$, на якій немає струмів провідності (рис.~\ref{pic:Refraction}). Нехай
силова лінія утворює кут $\alpha_1$ з нормаллю в середовищі 1 і кут $\alpha_2$ --- в середовищі 2. Граничні умови дають:
\begin{equation*}
	\begin{cases}
		H_{1\tau}=H_{2\tau} \ \Rightarrow\ H_1\sin\alpha_1 = H_2\sin\alpha_2, \\
		B_{1n}=B_{2n} \ \Rightarrow\ \mu_1H_1\cos\alpha_1 = \mu_2H_2\cos\alpha_2.
	\end{cases}
\end{equation*}

\begin{figure}[h!]\centering
	\localinput{refraction.tikz}
	\caption{Заломлення силових ліній $\Hfield$ і $\Bfield$ на межі розділу двох
		магнетиків\label{pic:Refraction}}
\end{figure}

Ділячи ці рівняння одне на одне, отримуємо \emphx{закон заломлення} силових ліній, однаковий і для $\Bfield$, і для $\Hfield$:
\begin{equation}\label{eq:RefractionLaw}
	\tcbhighmath{
		\frac{\tg\alpha_2}{\tg\alpha_1} = \frac{\mu_2}{\mu_1}.
	}
\end{equation}
У середовищі з більшим значенням $\mu$ силові лінії відхиляються на більший кут від нормалі до межі розділу.

\subsection*{Магнітне екранування}

На заломленні силових ліній ґрунтується явище \emphx{магнітного захисту (екранування)}. Якщо замкнену оболонку з феромагнетика (наприклад,
залізну) внести в зовнішнє магнітне поле, то, оскільки $\mu_\text{залізо} \gg \mu_\text{повітря}\approx 1$, за законом
заломлення~\eqref{eq:RefractionLaw} силові лінії всередині оболонки відхиляються майже до дотичної площини до її поверхні --- тобто
\emphx{\enquote{втягуються} в товщу оболонки} і концентруються переважно в ній самій, майже не проникаючи у внутрішню порожнину
(рис.~\ref{pic:MagShield}). У результаті поле всередині порожнини виявляється істотно ослабленим порівняно із зовнішнім.

\begin{SCfigure}[2][h!]\centering
	\includegraphics[width=0.25\linewidth]{MagShield}
	\caption{Магнітне екранування: силові лінії концентруються в товщі
		феромагнітної оболонки\label{pic:MagShield}}
\end{SCfigure}

\begin{Attention}
	Магнітне екранування --- принципово \emphx{інше} явище, ніж електростатичне екранування провідником: провідник в електростатиці повністю
	\emphx{виштовхує} поле зі свого об'єму (компенсуючи його вільними зарядами), тоді як феромагнітна оболонка, навпаки, \emphx{затягує} лінії
	поля в себе, немов пропонуючи їм \enquote{шлях найменшого опору} (за аналогією зі струмом, що воліє текти по провіднику з малим опором, а не
	по діелектрику). Саме тому магнітний захист вимагає \emphx{суцільної, замкненої} оболонки з високопроникного матеріалу --- на відміну від
	електростатичного захисту, для якого досить навіть тонкої провідної сітки; цей принцип широко застосовують для захисту чутливих приладів
	(осцилографів, детекторів, магнітометрів) від зовнішніх магнітних завад.
\end{Attention}



%\begin{Attention}
%	Феромагнітні властивості (і, зокрема, здатність до спонтанної намагніченості доменів) --- явище \emphx{не абсолютне}, а таке, що зникає з
%	підвищенням температури: для кожного феромагнетика існує \emphx{температура Кюрі} $T_C$ (для заліза $T_C\approx 770\,^\circ$C), вище якої
%	тепловий рух руйнує впорядкування магнітних моментів усередині доменів, і речовина перетворюється на звичайний \emphx{парамагнетик}. Це
%	прямий магнітний аналог температури Кюрі для сегнетоелектриків (гл.~\ref{FileldInDielectrics}), і, зрештою, наслідок тієї самої загальної
%	конкуренції між впорядковувальною взаємодією та тепловим хаосом.
%\end{Attention}

%%% --------------------------------------------------------
%\section{Підсумки розділу}
%%% --------------------------------------------------------
%
%\begin{table}[h!]\small\centering
%	\caption{Основні співвідношення теорії магнітного поля в речовині}
%		\begin{tblr}{
%			colspec={X[l] X[l]},
%			row{1} = {1.5em, c, azure2, fg=white,font=\bfseries\sffamily},
%			row{odd[2]}={gray9},
%			row{2-Z}={1.7em},
%			cell{1}{4}={mode=text},
%			hlines={azure3},
%			vlines={azure3},
%			vline{2-4}={1}{white},
%			hline{2}={white}
%		}
%		Вектор намагніченості                     & $\displaystyle \vect J = \frac1V\sum_i \vect p_i$                               \\
%		Циркуляція $\vect J$ (інтегральна форма)  & $\displaystyle c\oint\limits_L \vect J\cdot d\vect\ell = I_m$                    \\
%		Циркуляція $\vect J$ (диференціальна форма) & $\vect j_m = c\rot\vect J$                                                     \\
%		Вектор напруженості                       & $\Hfield = \Bfield - 4\pi\vect J$                                                \\
%		Для лінійних магнетиків                   & $\vect J = \chi_m\Hfield$                                                        \\
%		Зв'язок проникності й сприйнятливості      & $\mu = 1+4\pi\chi_m$                                                             \\
%		Для лінійних магнетиків                   & $\Bfield = \mu\Hfield$                                                           \\
%		Циркуляція $\Hfield$ у речовині (інт.)    & $\displaystyle \oint\limits_L \Hfield\cdot d\vect\ell = \frac{4\pi}c \iint\limits_S \vect j\cdot d\vect S$ \\
%		Циркуляція $\Hfield$ у речовині (диф.)    & $\rot\Hfield = \dfrac{4\pi}c\vect j$                                             \\
%		Гранична умова для $B_n$                  & $B_{1n}=B_{2n}$                                                                  \\
%		Гранична умова для $H_\tau$                & $H_{2\tau}-H_{1\tau} = \dfrac{4\pi}c i_b$                                        \\
%	\end{tblr}
%\end{table}

%% --------------------------------------------------------
\section{Мікроскопічна природа діа-, пара- і феромагнетизму}
%% --------------------------------------------------------

Досі ми користувалися сприйнятливістю $\chi_m$ як суто дослідною характеристикою речовини. Тепер, озброївшись формулами для магнітного
моменту й гіромагнітного відношення з гл.~\ref{MagFieldVacuum}, розкриємо обидва механізми кількісно, а наостанок з'ясуємо, звідки береться значно сильніше й принципово інше явище ---
спонтанне впорядкування у феромагнетиках.

\subsection*{Природа діамагнетизму: прецесія Лармора}

Розглянемо електрон, що рухається по орбіті навколо ядра; за формулою~\eqref{eq:GyromagneticClassical} (гл.~\ref{MagFieldVacuum}) його орбітальний
магнітний момент пов'язаний з орбітальним моментом імпульсу $\vect L$ гіромагнітним відношенням $\vect p = \gamma\vect L$, $\gamma=-e/(2m_ec)$. У
зовнішньому полі $\Bfield$ на цей момент діє момент сили~\eqref{eq:TorqueDipoleAnalogy}, який змінює $\vect L$ за рівнянням
\begin{equation*}
	\frac{d\vect L}{dt} = \vect p\times\Bfield = \gamma\left(\vect L\times\Bfield\right) = \vect\omega_L\times\vect L, \qquad
	\vect\omega_L\equiv-\gamma\Bfield.
\end{equation*}
Це --- рівняння рівномірної \emphx{прецесії} вектора $\vect L$ (а разом з ним --- і всієї площини орбіти) навколо напрямку $\Bfield$ з кутовою
швидкістю
\begin{equation}\label{eq:LarmorFreqMag}
	\tcbhighmath{
		\omega_L = |\gamma|B = \frac{eB}{2m_ec}
	}
\end{equation}
--- \emphx{ларморовою частотою}, однаковою для будь-якого електрона незалежно від радіуса й швидкості його орбіти (\emphx{теорема Лармора}).

%---------------------------------------------------------
\begin{figure}[h!]\centering
    \localinput{larmor_precession.tikz}
\caption{Прецесія орбіти електрона в зовнішньому полі $\Bfield$: магнітний
момент $\vect{p}_m$ (а разом з ним і площина орбіти) обертається навколо
$\Bfield$ з ларморовою частотою $\omega_L$. Ця прецесія еквівалентна
додатковому коловому струму $I=-e\omega_L/(2\pi)$, який створює індукований
магнітний момент $\Delta\vect p$, напрямлений проти $\Bfield$.}
\label{pic:LarmorPrecession}
\end{figure}
%---------------------------------------------------------


До ввімкнення поля електрон рухався по коловій орбіті; щойно поле накладено,
ця орбіта починає прецесувати (рис.~\ref{pic:LarmorPrecession}).  Прецесія орбіти --- це додатковий круговий рух заряду $-e$ з частотою $\omega_L$ навколо осі, паралельної $\Bfield$, тобто додатковий коловий
струм $I=-e\omega_L/(2\pi)$ на відстані $\rho$ від цієї осі. За формулою для моменту колового витка~\eqref{eq:MagMomentLoop} цей струм створює
додатковий (\emphx{індукований}) магнітний момент
\begin{equation}\label{eq:InducedDiaMag}
	\Delta\vect{p}_m = \frac{I\pi\rho^2}{c}\,\frac{\Bfield}{B} = -\frac{e\omega_L\rho^2}{2c}\,\frac{\Bfield}{B} = -\frac{e^2\rho^2}{4m_ec^2}\, \Bfield,
\end{equation}
напрямлений, зверніть увагу, \emphx{проти} поля $\Bfield$ --- незалежно від того, у який бік електрон обертався початково: це і є мікроскопічний
прояв правила Ленца.

\begin{Attention}
	Той самий результат~\eqref{eq:InducedDiaMag} випливає із закону
	електромагнітної індукції. При вмиканні поля потік крізь орбіту
	$\Phi=\pi r^2 B$ змінюється, і виникає ЕРС
	$\mathcal{E}=-(1/c)\,d\Phi/dt$. Вона створює дотичне поле
	$E_{\text{ін}}=r\dot B/(2c)$, яке за час $dt$ змінює швидкість електрона на
	\begin{equation*}
		dv = \frac{eE_{\text{ін}}}{m_e}\,dt
		= \frac{er}{2m_ec}\,dB .
	\end{equation*}
	Відповідна зміна струму $dI=-e\,dv/(2\pi r)$ дає
	\begin{equation*}
		d\Delta p = \frac{\pi r^2}{c}\,dI
		= -\frac{e^2 r^2}{4m_ec^2}\,dB
		\quad\Longrightarrow\quad
		\Delta p = -\frac{e^2 r^2}{4m_ec^2}\,B .
	\end{equation*}
	Знак мінус --- це правило Ленца: індукований момент напрямлений проти поля,
	яке його викликало.
\end{Attention}

Формула~\eqref{eq:InducedDiaMag} виведена для однієї класичної орбіти. У
квантовій механіці електрон не має траєкторії: його стан описується хвильовою
функцією, а заряд розподілений у просторі з густиною $-e|\psi|^2$. Тому
замість фіксованого $\rho$ слід підставити середнє $\langle\rho^2\rangle$ за
цією електронною хмарою. Електронна хмара сферично
симетрична, тож
\begin{equation*}
	\langle\rho^2\rangle = \tfrac23\langle r^2\rangle .
\end{equation*}
Підсумовуючи внесок усіх $Z$ електронів, отримуємо індукований момент атома,
а тоді й магнітну сприйнятливість речовини з концентрацією атомів $n$
(\emphx{формула Ланжевена для діамагнетизму}):
\begin{equation}\label{eq:LangevinDiaMag}
	\tcbhighmath{
		\chi_m^{\text{діа}} = -\frac{nZe^2\langle r^2\rangle}{6m_ec^2} < 0.
	}
\end{equation}

Ця величина не залежить від температури (розмір атома $\langle r^2\rangle$ від неї практично не залежить) і за порядком величини становить
$10^{-7}$--$10^{-6}$ (у гаусовій системі; у СІ відповідні значення в $4\pi$ разів більші) --- слабкий, але \emphx{універсальний} ефект,
наявний в атомах будь-якої речовини, включно з пара- і феромагнетиками.

%% --------------------------------------------------------
\subsection*{Природа парамагнетизму: закон Кюрі}
%% --------------------------------------------------------

Якщо атом має власний (нескомпенсований) магнітний момент $p_0\neq0$, до щойно знайденого діамагнітного внеску додається значно потужніший
ефект --- часткова \emphx{орієнтація} вже наявних моментів уздовж поля. Це --- дослівний магнітний аналог орієнтаційної поляризації полярних
діелектриків (гл.~\ref{FileldInDielectrics}, п.~\ref{sec:Langevin-Debye}): енергія моменту в полі дорівнює $U=-p_0B\cos\theta$
(формула~\eqref{eq:PotentialEnergyDipole}, гл.~\ref{MagFieldVacuum}), а рівноважний розподіл орієнтацій підпорядковується розподілу Больцмана.
У практично важливій межі слабкого поля й не надто низької температури ($p_0B\ll kT$) середня проекція моменту на напрямок поля лінійна за
полем, за тією самою викладкою, що й для електричного диполя~\eqref{eq:LangevinDebye}:
\begin{equation}\label{eq:ParaAvgMag}
	\langle p_z\rangle = \frac{p_0^2 B}{3kT}.
\end{equation}
Звідси намагніченість $n$ атомів в одиниці об'єму дорівнює $J=n\langle p_z\rangle$, а сприйнятливість спадає обернено пропорційно температурі ---
\emphx{закон Кюрі} (пряма аналогія закону~\eqref{eq:CurieLaw} для полярних діелектриків):
\begin{equation}\label{eq:CurieLawMag}
	\tcbhighmath{
		\chi_m^{\text{пара}} = \frac{np_0^2}{3kT} \equiv \frac{C}T > 0,
	}
\end{equation}
де $C=np_0^2/3k$ --- \emphx{стала Кюрі}. Оскільки за кімнатної температури $\chi_m^{\text{пара}}$, як правило, на один--два порядки перевищує
$\bigl|\chi_m^{\text{діа}}\bigr|$, саме вона визначає знак і величину сумарної сприйнятливості речовини, щойно в її атомах з'являється власний момент, --- а
слабший діамагнітний внесок~\eqref{eq:LangevinDiaMag} лишається практично непомітним на його тлі.

\subsection*{Природа феромагнетизму: обмінна взаємодія і поле Вейса}

Закон Кюрі~\eqref{eq:CurieLawMag} описує \emphx{незалежні} атомні моменти, що орієнтуються лише під дією зовнішнього поля й теплового хаосу. Але
у феромагнетику є локальні області --- \emphx{домени} --- типовим розміром від часток до сотень мікрометрів. У межах кожного домену магнітні
моменти всіх атомів спонтанно вишикувані в одному напрямку (рис.~\ref{pic:Domains}) --- без жодного зовнішнього поля. У різних доменах напрями
намагніченості, як правило, різні, тому сумарна намагніченість ненамагніченого зразка загалом близька до нуля. Доменна структура цілком
аналогічна до доменної структури сегнетоелектриків (гл.~\ref{FileldInDielectrics}).

\begin{figure}[h!]\centering
	\includegraphics[width=0.35\linewidth]{domains}
	\caption{Доменна структура феромагнетика: у кожному домені намагніченість стала, але напрями
		намагніченості різних доменів, як правило, різні\label{pic:Domains}}
\end{figure}


Формальне пояснення спонтанної орієнтації магнітних моментів у межах домену запропонував П.~Вейс: припустимо, що на кожен атомний момент,
окрім зовнішнього поля $\Hfield$, діє ще й
\emphx{молекулярне поле} $\Hfield_{\text{мол}} = \lambda\vect J$, пропорційне намагніченості \emphx{сусідніх} атомів (тут $\lambda$ --- стала
Вейса). Підставивши це ефективне поле $\Hfield+\lambda\vect J$ у закон Кюрі~\eqref{eq:CurieLawMag} замість $\Hfield$ (для слабких полів, коли
$B\approx H$), дістаємо самоузгоджене рівняння для намагніченості:
\begin{equation*}
	J = \frac CT\left(H+\lambda J\right) \quad\Longrightarrow\quad \chi_m = \frac JH = \frac{C}{T-C\lambda}.
\end{equation*}
Позначивши $T_C\equiv C\lambda$ --- \emphx{температуру Кюрі} (для заліза $T_C\approx770\,^\circ$C), отримуємо \emphx{закон Кюрі--Вейса}:
\begin{equation}\label{eq:CurieWeissLaw}
	\tcbhighmath{
		\chi_m = \frac{C}{T-T_C}.
	}
\end{equation}

Вище $T_C$ речовина парамагнітна: власні моменти атомів орієнтовані
хаотично, і без зовнішнього поля намагніченість дорівнює нулю. Коли
$T\to T_C^+$, сприйнятливість~\eqref{eq:CurieWeissLaw} прямує до
нескінченності --- парамагнітний стан втрачає стійкість. Щойно $T$
опускається нижче $T_C$, будь-яке мале відхилення $\vect J$ від нуля
підсилюється молекулярним полем швидше, ніж його руйнує тепловий рух, ---
і в системі самочинно встановлюється ненульова намагніченість навіть без
зовнішнього поля. Це --- феромагнітний стан.

Залишається питання: чим фізично є поле Вейса $\Hfield_{\text{мол}}$? Щоб упорядкування витримувало тепловий рух аж до кімнатної температури,
потрібне поле порядку $10^7$--$10^8$~Гс (з умови $p_0B\sim kT$). Натомість поле, яке реально створює сусідній атомний диполь на міжатомній
відстані, становить лише $\sim10^3$~Гс --- на чотири-п'ять порядків менше. Отже, \emphx{справжня} причина --- не магнітна
взаємодія диполів, а \emphx{обмінна взаємодія} --- суто квантовий ефект, що не має класичного аналога: принцип Паулі вимагає, щоб повна хвильова функція двох
електронів була антисиметричною, тому за \emphx{паралельних} спінів антисиметричною мусить бути її координатна частина, --- а це, своєю чергою, автоматично \emphx{розводить} такі електрони в просторі й тим самим \emphx{зменшує} енергію їхнього
кулонівського відштовхування порівняно з випадком антипаралельних спінів. У деяких кристалах (переважно в перехідних металах з незаповненою
$3d$-оболонкою --- \ce{Fe}, \ce{Co}, \ce{Ni}) цей суто електростатичний за походженням, але зумовлений квантовою статистикою виграш в енергії
виявляється достатньо великим, щоб \emphx{паралельна} орієнтація сусідніх атомних моментів була енергетично вигіднішою за хаотичну, --- саме це
і відіграє роль \enquote{молекулярного поля} Вейса.

\begin{Attention}
	Обмінна взаємодія --- \emphx{короткодійна} (діє практично лише між найближчими сусідами в ґратці) і в різних речовинах може сприяти як
	\emphx{паралельній}, так і \emphx{антипаралельній} орієнтації сусідніх моментів. Якщо вона вирівнює сусідні моменти паралельно --- маємо
	\emphx{феромагнетик}; якщо антипаралельно, причому моменти рівні за величиною, --- сумарна намагніченість зникає, і речовину називають
	\emphx{антиферомагнетиком}; якщо ж антипаралельно, але моменти різні за величиною (наприклад, належать різним типам вузлів ґратки), --- повної
	компенсації не відбувається, і виникає \emphx{феримагнетик} (класичний приклад --- магнетит $\ce{Fe3O4}$, --- та сама речовина, з якої
	виготовляли перші відомі людству магніти). У всіх трьох випадках --- та сама квантова природа взаємодії, лише різний \emphx{знак і
	геометрія} її прояву.
\end{Attention}



%\emphx{Феромагнетиками} називають тверді речовини, здатні мати \emphx{спонтанну намагніченість всередині окремих доменів} --- тобто намагнічені
%навіть за \emphx{відсутності} зовнішнього поля. Типові представники --- залізо, кобальт, нікель та багато їхніх сплавів.

\subsection*{Гістерезис у феромагнетиках}

\emphx{Гістерезисом} називають неоднозначну петлеподібну залежність намагніченості $J$ феромагнетика від зовнішнього поля $H$ за його циклічної
зміни (рис.~\ref{pic:Hysteresis}).

\begin{SCfigure}[1.25][h!]\centering
	\localinput{hysteresis.tikz}
	\caption{Петля гістерезису феромагнетика: $J_r$ --- залишкова намагніченість, $H_c$ ---
		коерцитивна сила\label{pic:Hysteresis}}
\end{SCfigure}

Проходження петлі відбувається за такими стадіями:
\begin{enumerate}
	\item За достатньо великого поля $H$ намагніченість досягає \emphx{насичення}: усі домени вже зорієнтовані вздовж поля, і зростання $J$
	      практично припиняється, $J\to J_s=\const$. Індукція $\Bfield = \Hfield+4\pi\vect J$ при цьому й далі зростає --- але вже лише за
	      рахунок самого $\Hfield$, з таким самим нахилом, як у вакуумі.
	\item Коли поле зменшують до нуля ($H=0$), намагніченість \emphx{не} повертається до нуля, а залишається на рівні залишкової намагніченості
	      $J_r$ --- саме тому постійні магніти взагалі можуть існувати.
	\item Щоб звести намагніченість до нуля, потрібно прикласти поле протилежного напрямку --- \emphx{коерцитивну силу} $-H_c$.
	\item З подальшим зростанням поля в тому самому (від'ємному) напрямку намагніченість знову зростає за модулем до насичення $-J_s$.
	\item Коли від'ємне поле зменшують за модулем до нуля, намагніченість знову залишається ненульовою, на рівні $-J_r$, --- і цикл, після проходження
	      додатної коерцитивної сили $+H_c$, замикається.
\end{enumerate}

\begin{Attention}
	Площа петлі гістерезису має пряму фізичну інтерпретацію: вона дорівнює (з точністю до множника) енергії, що незворотно розсіюється у вигляді
	тепла за один повний цикл перемагнічування --- саме тому вузька петля означає малі втрати, а широка --- великі.
\end{Attention}

\subsection*{Магнітом'які та магнітотверді феромагнетики}

Залежно від ширини петлі гістерезису феромагнетики поділяють на дві категорії.

\emphx{Магнітом'які} матеріали (рис.~\ref{pic:Soft}), які застосовують як осердя в трансформаторах, електромагнітах та інших приладах, мають дуже
\emphx{вузьку} петлю гістерезису. Вузькість петлі (мала коерцитивна сила $H_c$) означає малі втрати на перемагнічування --- важливо, наприклад,
для трансформаторів, які перемагнічуються десятки чи сотні разів за секунду. На ділянці до насичення петля близька до прямої лінії, для якої, як і
для неферомагнітних речовин, можна писати $\Bfield = \mu\Hfield$.

\begin{figure}[h!]
	\centering
	\begin{subfigure}[t]{0.48\linewidth}
		\centering
		\includegraphics[width=0.7\linewidth]{Soft}
		\caption{Вузька петля гістерезису магнітом'якого матеріалу}
		\label{pic:Soft}
	\end{subfigure}\hfill
	\begin{subfigure}[t]{0.48\linewidth}
		\centering
		\includegraphics[width=0.7\linewidth]{Hard}
		\caption{Широка петля гістерезису магнітотвердого матеріалу}
		\label{pic:Hard}
	\end{subfigure}
	\caption{Петлі гістерезису магнітом'якого та магнітотвердого матеріалів}
	\label{pic:HysteresisComparison}
\end{figure}

Постійні ж магніти, навпаки, виготовляють із \emphx{магнітотвердих} матеріалів (рис.~\ref{pic:Hard}), наприклад із загартованої сталі, з
широкою петлею гістерезису й великою коерцитивною силою $H_c$ --- саме це забезпечує збереження намагніченості магніту з часом і під дією
випадкових зовнішніх завад.

%\begin{Attention}
%	Підсумуємо принципову відмінність трьох явищ. \emphx{Діамагнетизм} --- єдиний ефект, що не потребує жодного власного моменту: зовнішнє поле
%	саме \emphx{індукує} момент (через ларморову прецесію орбіт), завжди спрямований проти поля, --- тому він слабкий, від'ємний, універсальний і
%	не залежить від температури. \emphx{Парамагнетизм} потребує вже наявного власного моменту, який поле лише \emphx{частково орієнтує},
%	долаючи тепловий хаос, --- тому він на порядки сильніший за діамагнетизм, додатний і спадає з температурою (закон Кюрі). \emphx{Феромагнетизм}
%	--- це парамагнетизм, \emphx{підсилений} внутрішньою квантовою (обмінною) взаємодією настільки, що впорядкування настає \emphx{самочинно},
%	без зовнішнього поля, аж до значної (для заліза --- майже 800\textdegree{}C) температури Кюрі; вище неї обмінна взаємодія програє тепловому
%	руху, і речовина повертається до звичайної парамагнітної поведінки за законом Кюрі--Вейса~\eqref{eq:CurieWeissLaw}.
%\end{Attention}

%% --------------------------------------------------------
\section*{Підсумки}
%% --------------------------------------------------------

Уся теорія магнітного поля в речовині побудована за єдиним лекалом, яке ми вже бачили в теорії діелектриків: (1)~ввести густину
мікроскопічних джерел поля (там --- $\vect P$, тут --- $\vect J$); (2)~виразити через неї зв'язані (молекулярні) джерела --- там через
дивергенцію ($\rho' = -\div\vect P$), тут через ротор ($\vect j_m = c\rot\vect J$) --- відмінність, що відображає потенціальну природу $\Efield$
проти вихрової природи $\Bfield$; (3)~ввести допоміжний вектор ($\Dfield$ або $\Hfield$), теорема для якого містить \emphx{лише} вільні
(вимірювані) джерела; (4)~для лінійних середовищ пов'язати допоміжний і основний вектори простим коефіцієнтом ($\epsilon$ або $\mu$). Розуміння
цієї спільної структури --- найкращий спосіб не заплутатися в позначеннях і знаках обох розділів курсу.